%%%%%%%%%%%%%%%%%%%%%%%%%%%%%%%%%%%%%%%%%%%%%%%%%%%%%%%%%%%%%%%%%%%%%%%%%%%%%%%%
%%  ZWISCHENABGABE: Einlage: Eine wissenschaftliche Arbeit   (Moodle: 11.11.2026)
%%  Aufbau nach pp-portfolio-geruest-einlage-2026-2027.tex.
%%  Der Text steht in einlagen/2_wiss_arbeit.tex und wird 1:1 in die Endabgabe übernommen.
%%  Overleaf: Menü -> Main document = diese Datei. PDF für Moodle benennen:
%%  Manevski_Wiss_Arbeit.pdf
%%%%%%%%%%%%%%%%%%%%%%%%%%%%%%%%%%%%%%%%%%%%%%%%%%%%%%%%%%%%%%%%%%%%%%%%%%%%%%%%
\newcommand{\dasDatum}{11.\,11.\,2026}   % Moodle-Frist dieser Abgabe
%%%%%%%%%%%%%%%%%%%%%%%%%%%%%%%%%%%%%%%%%%%%%%%%%%%%%%%%%%%%%%%%%%%%%%%%%%%%%%%%
%%  EINSTELLUNGEN – die einzige Datei, in der du deine Angaben einträgst.
%%  Wird von ALLEN Hauptdateien (Endabgabe + Zwischenabgaben) eingelesen.
%%  Entspricht dem Block "Die Angaben" der offiziellen Gerüste
%%  pp-portfolio-geruest-2026-2027.tex / pp-portfolio-geruest-einlage-2026-2027.tex
%%%%%%%%%%%%%%%%%%%%%%%%%%%%%%%%%%%%%%%%%%%%%%%%%%%%%%%%%%%%%%%%%%%%%%%%%%%%%%%%

%%%%%%%%%%%%%%%%%%%%%%%%%%%%%%%%%%%%% Die Angaben %%%%%%%%%%%%%%%%%%%%%%%%%%%%%%%
\newcommand{\Semester}{WH 2026-2027}
\newcommand{\Projekt}{Dynamik Keltischer Wackelsteine}   % Titel laut Themenliste WS 26/27
%% Schwerpunkt abgeleitet aus der Aufgabenverteilung im Arbeitsplan (1.4 Bewegungsgleichungen,
%% 1.6 Simulation Umkehrdynamik, 2.1.3/2.1.4 Firmware + Kalibrierung, 5.1 Datenpipeline) – prüfen!
\newcommand{\Schwerpunkt}{Bewegungsgleichungen und numerische Simulation der Spin-Umkehr}
\newcommand{\SchwerpunktKurz}{Simulation der Spin-Umkehr}        % kurze Form für die Kopfzeile
\newcommand{\IhrName}{Kristijan Manevski}                         % Vorname Nachname
\newcommand{\Matrikelnummer}{\platzhalter{Matrikelnummer}}
\newcommand{\Studiengang}{\platzhalter{Studiengang + Semester}}   % z.B. Physik B.Sc., 5. Semester
\newcommand{\emailAdresse}{\platzhalter{Uni-E-Mail-Adresse (\dots @uni-ulm.de)}}
%% Abgabedatum: jede Zwischenabgabe-Hauptdatei setzt ihr eigenes \dasDatum (Moodle-Frist);
%% ohne Vorgabe gilt das Datum der Endabgabe.
\providecommand{\dasDatum}{15.\,03.\,2027}

%% Weitere Angaben, die in mehreren Einlagen gebraucht werden
\newcommand{\Partner}{Janosch Lustig}                             % für Abgrenzung der Schwerpunkte
%% laut Titelseite von Janoschs Arbeitsplan-Abgabe (28.10.2026)
\newcommand{\PartnerSchwerpunkt}{Umkehrzeit und Parameterabhängigkeit}
\newcommand{\Laborbuchlink}{https://link-zum-elektronischen-laborjournal}  % Pflicht in Einlage 3!
%% Einlage 3: ein Praktikumstag (Do 10--12 Uhr, O26/233 bzw. O26/234), ≠ Tag von Janosch.
%% Termine: 15.10., 22.10., 29.10., 05.11., 12.11., 19.11., 26.11., 03.12., 10.12., 17.12.2026,
%%          14.01., 21.01., 28.01., 04.02., 11.02.2027 (Abgabe Protokoll: 09.12.2026)
\newcommand{\Protokolldatum}{TT.MM.2026}
\newcommand{\PaperKey}{KondoNakanishi2017}      % Einlage 2: Kondo & Nakanishi, Phys. Rev. E 95, 062207
\newcommand{\Vortragsdatum}{28.\,01.\,2027}      % Einlage 5
\newcommand{\FremdVortragKey}{VORLAGE-Vortrag}  % Einlage 6: bib-Key des anderen Vortrags

%%%%%%%%%%%%%%%%%%%%%%%%%%%%%%%%%%%%%%%%%%%%%%%%%%%%%%%%%%%%%%%%%%%%%%%%%%%%%%%%
%%  MODUS-SCHALTER
%%
%%  \Abgabefalse  = Entwurf: grüne Hilfen sichtbar, Kontrollseite mit
%%                  Umfangsmessung, Platzhalter rot markiert.
%%  \Abgabetrue   = Abgabe: grüne Farbe -> schwarz (wie offizielle Zeile 15/16),
%%                  alle Hilfen + Kontrollseite + Umfangsboxen verschwinden.
%%                  Übrig gebliebene \platzhalter{} bleiben ROT sichtbar und
%%                  erzeugen eine Warnung -> dann ist noch etwas offen!
%%
%%  \Hinweisefalse blendet im Entwurf nur die grünen Hilfstexte aus
%%                  (nützlich, um den Umfang realistisch zu messen).
%%%%%%%%%%%%%%%%%%%%%%%%%%%%%%%%%%%%%%%%%%%%%%%%%%%%%%%%%%%%%%%%%%%%%%%%%%%%%%%%
\newif\ifAbgabe
\newif\ifHinweise

\Abgabefalse      % <-- für die Abgabe auf \Abgabetrue setzen
\Hinweisetrue     % <-- Hilfstexte im Entwurf ein-/ausblenden

\newcommand{\dieEinlage}{Einlage: Eine wissenschaftliche Arbeit}

\documentclass[pdftex,12pt,twoside,footsepline,headsepline,a4paper,final]{scrreprt}
%%%%%%%%%%%%%%%%%%%%%%%%%%%%%%%%%%%%%%%%%%%%%%%%%%%%%%%%%%%%%%%%%%%%%%%%%%%%%%%%
%%  PRÄAMBEL – gemeinsam für Endabgabe und alle Zwischenabgaben.
%%  Teil 1 ist 1:1 die Präambel der offiziellen Gerüste (nicht verändern),
%%  Teil 2 ergänzt nur Pakete für Gleichungen, Bilder, PDF-Einbindung usw.,
%%  Teil 3 enthält die Hilfsmakros dieser Vorlage.
%%  Die \documentclass-Zeile steht (unverändert offiziell) in jeder Hauptdatei.
%%%%%%%%%%%%%%%%%%%%%%%%%%%%%%%%%%%%%%%%%%%%%%%%%%%%%%%%%%%%%%%%%%%%%%%%%%%%%%%%

%%%%%%%%%%%%%%%%%%%%%%%%%%%%% Teil 1: offizielle Präambel %%%%%%%%%%%%%%%%%%%%%%
\usepackage[verbose,a4paper,bottom=25mm,top=25mm]{geometry} %showframe,
\usepackage{scrlayer-scrpage}
\usepackage[T1]{fontenc}
\usepackage[utf8]{inputenc}
\usepackage{ngerman}
\usepackage{color}
\usepackage{amsfonts}
\usepackage{amssymb}
\usepackage{amsthm}
%%%%%%%%%%%%%%%%%%%%%%%%% Biblatex (offizielle Optionen, Backend Biber) %%%%%%%%%%
%% Overleaf erkennt biber automatisch. Zitierstil alphabetic: [Sch02]
\usepackage[sorting=nyt, backend=biber, citestyle=alphabetic, style=alphabetic,
            autocite=inline, labelalpha=true, hyperref=auto, backref=true,
            language=ngerman]{biblatex}
\addbibresource{literatur.bib}

%%%%%%%%%%%%%%%%%%%%%%%%%%%%% Teil 2: Zusatzpakete %%%%%%%%%%%%%%%%%%%%%%%%%%%%%%
\usepackage{amsmath}               % \eqref, align, nummerierte Gleichungen
\usepackage{xcolor}
\definecolor{darkgreen}{rgb}{0,0.5,0}   % offizielle Hinweisfarbe
\usepackage{graphicx}              % \includegraphics
\graphicspath{{abbildungen/}}
\usepackage{booktabs}              % schöne Tabellen
\usepackage[locale=DE]{siunitx}    % \SI{3.4}{\kilo\hertz}, \num{1.2e3}
\usepackage{pdfpages}              % \includepdf für Anhänge (Paper, Laborjournal, Folien)
\usepackage{pdflscape}             % Querformat: \begin{landscape}...\end{landscape}
\usepackage{pgfgantt}              % Gantt-Diagramm für den Arbeitsplan
\usepackage{environ}
\usepackage[savepos,abspage]{zref}
\usepackage[hidelinks]{hyperref}   % klickbare Verweise; zuletzt laden

%% Zitierregel (Marti): bei Journalartikeln DOI angeben, aber keine Webadresse.
%% -> url/urldate werden bei @article automatisch unterdrückt.
\AtEveryBibitem{%
  \ifentrytype{article}{\clearfield{url}\clearfield{urlyear}%
    \clearfield{urlmonth}\clearfield{urlday}}{}}

%% Kopf- und Fußzeile wie offizielles Gerüst (\ohead wird in der Hauptdatei gesetzt)
\pagestyle{scrheadings}
\ihead{\thepage}
\ofoot{\pagemark}
\ifoot{\zubearbeiten{\IhrName}}

%%%%%%%%%%%%%%%%%%%%%%%%%%%%% Teil 3: Hilfsmakros %%%%%%%%%%%%%%%%%%%%%%%%%%%%%%%
\makeatletter

%% --- offizieller Farbschalter (Zeile 12-15 der Gerüste), automatisch per Modus
\ifAbgabe
  \Hinweisefalse
  \newcommand{\zubearbeiten}[1]{\textcolor{black}{#1}}
\else
  \newcommand{\zubearbeiten}[1]{\textcolor{darkgreen}{#1}}
\fi

%% --- Grüne Hilfstexte: verschwinden im Abgabemodus vollständig
%%     \begin{hinweise} ... \end{hinweise}   (darf itemize enthalten)
%%     \hinweis{kurzer Text}
\NewEnviron{hinweise}{\ifHinweise\par{\color{darkgreen}\small\BODY\par}\fi}
\newcommand{\hinweis}[1]{\ifHinweise\par{\color{darkgreen}\small #1\par}\fi}

%% --- Platzhalter für eigenen Text: bleibt IMMER rot sichtbar (auch im
%%     Abgabemodus), damit nichts Unfertiges abgegeben wird.
\newcounter{platzhalter}
\DeclareRobustCommand{\platzhalter}[1]{%
  \stepcounter{platzhalter}%
  \textcolor{red}{\textbf{[}\,#1\,\textbf{]}}%
  \ifAbgabe\PackageWarning{Portfolio}{Platzhalter noch vorhanden: #1}\fi}
\AtEndDocument{%
  \ifnum\value{platzhalter}>0
    \PackageWarningNoLine{Portfolio}{Es sind noch \theplatzhalter\space
      Platzhalter im Dokument!}%
  \fi
  \immediate\write\@auxout{\string\gdef\string\pf@platzhalteralt{\theplatzhalter}}}
\providecommand{\pf@platzhalteralt}{?}

%% --- Umfangsmessung
%%   \begin{umfang}{Bezeichnung}{min}{max} ... \end{umfang}
%%   misst die Länge des eingeschlossenen Textes in Seiten (über die
%%   Position auf der Seite, auf 0.05 Seiten genau) und zeigt sie im Entwurf
%%   in einem Kasten an. Achtung: Abbildungen/Tabellen zählen nach den Regeln
%%   NICHT zur Länge, werden hier aber mitgemessen -> dann abziehen.
\zref@addprop{savepos}{abspage}
\newcounter{pf@umfang}
\newcommand*{\pf@umfangliste}{}
\newcommand*{\pf@umfangeintrag}[4]{% Name, min, max, ist
  \g@addto@macro\pf@umfangliste{\pf@umfangzeile{#1}{#2}{#3}{#4}}}
\newcommand*{\pf@setfarbe}[3]{% min max ist -> \pf@farbe
  \def\pf@farbe{darkgreen}%
  \ifdim #3pt<#1pt \def\pf@farbe{red}\fi
  \ifdim #3pt>#2pt \def\pf@farbe{red}\fi}
\newcommand*{\pf@umfangzeile}[4]{%
  #1 & #2 -- #3 & \pf@setfarbe{#2}{#3}{#4}\textcolor{\pf@farbe}{\textbf{#4}}\\}
\newenvironment{umfang}[3]{%
  \stepcounter{pf@umfang}%
  \edef\pf@uid{pfumf\the\value{pf@umfang}}%
  \def\pf@uname{#1}\def\pf@umin{#2}\def\pf@umax{#3}%
  \par\zsavepos{\pf@uid a}\nobreak
}{%
  \par\zsavepos{\pf@uid e}%
  \ifAbgabe\else
    \zref@refused{\pf@uid a}\zref@refused{\pf@uid e}%
    \edef\pf@ist{\fpeval{round(
        (\zref@extractdefault{\pf@uid e}{abspage}{0}
        -\zref@extractdefault{\pf@uid a}{abspage}{0})
        +(\zref@extractdefault{\pf@uid a}{posy}{0}
        -\zref@extractdefault{\pf@uid e}{posy}{0})/\number\textheight ,2)}}%
    \protected@write\@auxout{}{\string\pf@umfangeintrag{\pf@uname}{\pf@umin}{\pf@umax}{\pf@ist}}%
    \par\smallskip\noindent
    \pf@setfarbe{\pf@umin}{\pf@umax}{\pf@ist}%
    \fcolorbox{\pf@farbe}{white}{\scriptsize
      Umfang \emph{\pf@uname}: $\approx$\,\pf@ist~Seiten
      (Soll: \pf@umin\,--\,\pf@umax~Seiten, ohne Titel/Abb./Tab./Legenden)}%
    \par\smallskip
  \fi}

%% --- Kontrollseite (nur im Entwurf; ersetzt das grüne Kapitel
%%     "Hinweise zur Arbeit mit dieser Datei", das abgegeben nicht vorkommen darf)
\newcommand{\kontrollseite}{%
  \ifAbgabe\else
  \chapter*{Kontrollseite (erscheint nur im Entwurf)}
  {\color{darkgreen}
  Modus: \textbf{Entwurf}. Für die Abgabe in \texttt{einstellungen.tex}
  \texttt{\textbackslash Abgabetrue} setzen. Dann wird alles Grüne schwarz bzw.\
  entfernt (Regel: \enquote{Abgaben mit grüner Farbe oder stehen gebliebenen
  Hilfen ergeben einen Abzug}).\par\medskip
  \textbf{Offene Platzhalter} (letzter Durchlauf): \pf@platzhalteralt\par\medskip
  \textbf{Umfangskontrolle} (Werte aus dem letzten Kompilierdurchlauf;
  Hilfstexte stehen außerhalb der Messung):\par\smallskip
  \ifx\pf@umfangliste\@empty
    \emph{noch keine Messwerte -- nochmals kompilieren.}
  \else
    {\small\begin{tabular}{p{0.55\textwidth}cc}
      \toprule Teil & Soll [Seiten] & Ist [Seiten]\\\midrule
      \pf@umfangliste
      \bottomrule
    \end{tabular}}
  \fi\par\medskip
  %% Checkliste – erscheint nur auf der Kontrollseite im Entwurf.
%% Quellen: Portfolios im Projektpraktikum 2026-2027 (Marti, 09.07.2026),
%% offizielle Gerüste 2026-2027, Zitierregeln (31.08.2026),
%% Abbildungen und Gleichungen (Limmer 2007), Moodle-Terminliste.

\textbf{Fristen (Moodle; nächste Abgabe wird erst nach der vorherigen freigeschaltet!)}\par\smallskip
{\small\begin{tabular}{@{}ll@{}}
21.\,10.\,2026 & Motivationsformular (ohne Reflexion) -- obligatorisch\\
28.\,10.\,2026 & Arbeitsplan + Reflexion -- obligatorisch\\
11.\,11.\,2026 & Wissenschaftliche Arbeit + Reflexion -- obligatorisch\\
09.\,12.\,2026 & Protokoll eines Praktikumstages + Reflexion -- obligatorisch\\
20.\,01.\,2027 & Messergebnis + Reflexion -- obligatorisch\\
23.\,01.\,2027 & Späteste Anmeldung mündliche Leistung (Portal)\\
27.\,01.\,2027 & Folien des Vortrags + Materialien für Mitstudierende\\
28.\,01.\,2027 & Vorträge (Prüfung, Anwesenheitspflicht)\\
07.\,02.\,2027 & Folien mit Kommentar + Reflexion (vorgeschlagen, freiwillig)\\
21.\,02.\,2027 & Aussage aus anderem Vortrag + Reflexion (vorgeschlagen, freiwillig)\\
10.\,03.\,2027 & Späteste Anmeldung schriftliche Leistung (Portal)\\
\textbf{15.\,03.\,2027} & \textbf{Endabgabe Portfolio -- verspätet = nicht bestanden}\\
\end{tabular}}\par\smallskip
{\small Praktikumstage (Google Kalender): donnerstags 10--12\,Uhr, O26/233 bzw.\ O26/234 --
15.\,10.\ bis 17.\,12.\,2026 und 14.\,01.\ bis 11.\,02.\,2027; keine Termine 18.\,12.--13.\,01.
Die Anmeldetermine im Hochschulportal (23.\,01./10.\,03.) stammen aus dem Plan von Petros -- im Portal prüfen!}
\par\medskip

\textbf{Vor jeder Abgabe prüfen}
\begin{itemize}\itemsep0pt
\item \texttt{\textbackslash Abgabetrue} gesetzt, keine grüne Farbe, keine Hilfen, \emph{0 Platzhalter}.
\item Deckblatt vollständig: Name, Vorname, Matrikelnummer, Studiengang, E-Mail, Thema (Projekt + Schwerpunkt).
\item Inhaltsverzeichnis mit Seitenzahlen, auch der Anhänge; durchgehend Seitenzahlen.
\item Längen eingehalten (siehe Tabelle oben; Abbildungen/Tabellen/Legenden zählen nicht mit):
      Einleitung 2--4, Beschreibung $\tfrac14$--$\tfrac12$, Diskussion der Physik 1--3,
      Reflexion 1--2, Gesamtreflexion 3--6 Seiten. Schriftgröße 12\,pt (fest eingestellt).
\item Beschreibung und Reflexion sind \emph{keine Nacherzählung}; die Leitfragen werden in Fließtext
      beantwortet (nicht in Einzelsätzen), erkennbar zugeordnet.
\item Paper (Einlage 2) wird in den späteren Einlagen zur Argumentation verwendet; Gleichungen werden
      zur Argumentation genutzt.
\item Peer-Feedback (Kommilitonen) eingeholt, aufgegriffen -- und zitiert. Petros' Arbeitsplan als Vorlage
      zitieren (Key \texttt{Petros2026Arbeitsplan}); Schwerpunkt-Abgrenzung mit Janosch festhalten.
\item Jede Übernahme zitiert: Literatur, Bilder (\texttt{\textbackslash cite[Abb.~3.4]\{key\}}),
      Gleichungen (\texttt{\textbackslash cite[Gl.~(3)]\{key\}}), Aussagen von Personen, Anleitungen,
      Datenblätter, Webseiten (mit Abrufdatum), \textbf{jede KI-Nutzung mit System und ALLEN Prompts}.
\item Journalartikel: Autoren, Titel, Journal, Band, Seiten, Jahr, DOI (nicht als Link).
\item Abbildungen: Bildunterschrift, im Text referenziert (\texttt{\textbackslash ref}), Schrift
      $\geq 10$\,pt, Achsen mit Einheit als \enquote{$f$ / kHz} (nicht \enquote{[kHz]}), ggf.\ Residuenplot.
\item Tabellen: Überschrift, im Text referenziert; Messdaten nicht als lange Listen im Text, sondern im Anhang.
\item Wichtige Gleichungen nummeriert, alle Symbole definiert und einheitlich verwendet.
\item Dateiname für Moodle: \texttt{Manevski\_WelcheEinlage.pdf} (nur PDF oder DOCX).
\item Bis auf den Arbeitsplan müssen sich alle Einlagen von denen der Partnerin/des Partners unterscheiden.
\end{itemize}

  }%
  \fi}
\providecommand{\enquote}[1]{\glqq #1\grqq}

%% --- Anhang aus PDF mit durchgehenden Seitenzahlen (Regel: "Die Arbeit ist
%%     durchgängig mit Seitenzahlen versehen ... auch der Anhänge")
%%   \anhangpdf{Kapiteltitel}{label}{datei.pdf}           alle Seiten
%%   \anhangpdf[pages={1-3}]{Kapiteltitel}{label}{datei.pdf}
%%   Die Kapitelüberschrift steht auf einer eigenen Seite, danach folgen die
%%   PDF-Seiten (auf 82 % verkleinert, damit Kopf-/Fußzeile frei bleiben).
\newcommand{\anhangpdf}[4][pages=-]{%
  \chapter{#2}\label{#3}%
  \includepdf[#1,scale=0.82,pagecommand={\thispagestyle{scrheadings}}]{#4}}
%% Gleiches ohne neues Kapitel (für weitere PDFs im selben Anhang)
\newcommand{\pdfseiten}[2][pages=-]{%
  \includepdf[#1,scale=0.82,pagecommand={\thispagestyle{scrheadings}}]{#2}}

\makeatother

%% Antwortfeld im Motivationsformular (Anhang der Endabgabe)
\newcommand{\antwort}[1]{\par\smallskip #1\par\bigskip}

\ohead{\zubearbeiten{\Projekt{} - \SchwerpunktKurz}}

\begin{document}

%% Titelseite einer ZWISCHENABGABE – wörtlich aus pp-portfolio-geruest-einlage-2026-2027.tex
%% \dieEinlage wird in der jeweiligen Hauptdatei (Manevski_....tex) gesetzt.
\title{Universit\"{a}t Ulm\\\vspace{1cm}
{\LARGE \textnormal{\textsc{Fakult\"{a}t f\"{u}r Naturwissenschaften}}}\\\vspace{2cm} {\LARGE \textnormal{\textsl{Projektpraktikum \Semester}}}\\\vspace{2cm} {\LARGE\zubearbeiten{\textit{\Projekt}}}\\
{\Large \zubearbeiten{\textit{\Schwerpunkt}}}\\\vspace{0.5cm} \zubearbeiten{\textit{\dieEinlage}}}

\author{{\LARGE \zubearbeiten{\IhrName}}\\ $_{}$\\$_{}$\\\vspace*{1.cm}
Matrikelnummer:  \zubearbeiten{\Matrikelnummer}\\\vspace{0.8cm} Studiengang: \zubearbeiten{\Studiengang}\\\vspace{0.4cm} e-Mail: \zubearbeiten{\emailAdresse}}

\date{\zubearbeiten{\dasDatum}}

\maketitle

\tableofcontents
\kontrollseite   % nur im Entwurf sichtbar

\chapter{Hauptteil}
\begin{hinweise}
Beachten Sie die Hinweise in der Datei mit dem Gerüst für die Endabgabe! Auch bei einer Zwischenabgabe ist
eine korrekte Literaturarbeit notwendig. Denken Sie daran, dass Sie alle Hinweise im Text auf andere Werke
(Artikel, Bücher, Bilder, Grafiken) aber auch alle Erwähnungen von Meinungen anderer Personen mit einem Zitat
versehen müssen! Beachten Sie auch die Fragen bei der jeweiligen Abgabe in Moodle.
\end{hinweise}

%%%%%%%%%%%%%%%%%%%%%%%%%%%%%%%%%%%%%%%%%%%%%%%%%%%%%%%%%%%%%%%%%%%%%%%%%%%%%%%%
%%  EINLAGE 2: WISSENSCHAFTLICHE ARBEIT     Zwischenabgabe: 11.11.2026 (obligatorisch)
%%  Struktur: Beschreibung (1/4–1/2 S.) + Diskussion der Physik (1–3 S.) + Reflexion (1–2 S.)
%%  Paper: Kondo & Nakanishi, Phys. Rev. E 95, 062207 (2017) -> \PaperKey = KondoNakanishi2017.
%%  Janosch darf NICHT dasselbe Paper nehmen. Falls doch: auf Garcia & Hubbard (1988),
%%  Key GarciaHubbard1988, wechseln (liegt ebenfalls in Literatur/).
%%%%%%%%%%%%%%%%%%%%%%%%%%%%%%%%%%%%%%%%%%%%%%%%%%%%%%%%%%%%%%%%%%%%%%%%%%%%%%%%
% \PaperKey wird in einstellungen.tex gesetzt

\section{Einlage: Eine wissenschaftliche Arbeit zu einem Teilaspekt Ihres Schwerpunktgebietes}
\label{sec:paper}

\begin{hinweise}
Diese Einlage muss aus einer begutachteten Fachzeitschrift sein. Sie sollen mit dieser Einlage die
professionelle Literaturrecherche lernen. Bücher, Konferenzbeiträge, Trade Journals und ähnliches werden hier
nicht akzeptiert.
Dieser Artikel \textbf{muss} in Ihrer weiteren Arbeit verwendet werden (v.\,a. in Einlage 4: \enquote{Eine
Bezug zu Ihrem Paper ist Voraussetzung}).
Sie dürfen nicht den gleichen Artikel verwenden wie Ihre Partnerin / Ihr Partner.
Akzeptable wissenschaftliche Einlagen sind begutachtete wissenschaftliche Artikel aus wissenschaftlichen
Zeitschriften, die Sie über webofknowledge.com oder auch über Google Scholar finden. Erfüllt die Einlage die
Anforderung nicht, wird sie in der Bewertung nicht berücksichtigt. Bei Zweifeln über die Zulässigkeit der
beiden Einlagen mit wissenschaftlichen Arbeiten fragen Sie Prof.~O.~Marti.
\end{hinweise}

\subsection{Beschreibung}
\begin{hinweise}
Die Beschreibung gibt alle notwendigen Informationen an, die die Einlage beschreiben, also zum Beispiel den
Fundort, die Literaturstelle und eine kurze Inhaltsübersicht, aber keine Nacherzählung ($1/4$ Seite, max.
$1/2$ Seite!). Leitfaden: Art der Einlage, wie bin ich darauf aufmerksam geworden (Datenbank, Suchbegriffe,
Hinweis von \dots), wesentliche Aussagen in eigenen Worten.
\end{hinweise}
\begin{umfang}{Einlage 2 (Paper): Beschreibung}{0.25}{0.5}
\platzhalter{In eigenen Worten ausformulieren -- Stichpunkte: Y.~Kondo und H.~Nakanishi,
\enquote{Rattleback dynamics and its reversal time of rotation}, Phys.\ Rev.\ E \textbf{95}, 062207
(2017) \cite{\PaperKey}, begutachtete Fachzeitschrift (APS); Fundort: Literaturrecherche in Web of
Science/Google Scholar (Arbeitspaket 1.1), Suchbegriffe: \dots; Inhalt: physikalisch transparente
Neuformulierung der Rattleback-Dynamik, Reduktion auf drei Variablen (Drehen, Nicken, Rollen),
Garcia-Hubbard-Formel für die Umkehrzeit als Produkt von vier Faktoren, ausführliche numerische
Simulationen zu Gültigkeit und Grenzen der Formel; PDF in Anhang~\ref{anh:paper}}
\end{umfang}

\subsection{Diskussion der Physik}
\begin{hinweise}
Beschreiben Sie hier die Physik dieser Einlage. Fassen Sie den Artikel zusammen. Arbeiten Sie besonders die
Physik, die für ihr Schwerpunktthema relavant ist, heraus.
\begin{itemize}\itemsep0pt
\item Was ist die Hauptaussage der wissenschaftlichen Arbeit?
\item Was sind die Voraussetzungen für den Inhalt des Artikels?
\item Wie hilft physikalisch die Physik dieser Einlage bei Ihrem Schwerpunktthema?
\item Wo gab es Schwierigkeiten beim Verständnis der Physik? Welche? Mit welcher zusätzlichen Literatur haben Sie das Problem gelöst
\item Haben Sie weitere Artikel oder Bücher konsultiert um das Paper zu verstehen? Welche?
\item Welche Aussagen verwenden Sie für Ihr Schwerpunktthema und wo? Warum?
\end{itemize}
\textbf{Die Diskussion der Physik sollte mindestens eine Seite und maximal 3 Seiten lang sein. Die
Längenangaben sind ohne Titel, Bilder, Tabellen und Tabellen- oder Bildlegenden und ohne übergrosse Abstände
zu Gleichungen zu verstehen.}
Gleichungen zur Argumentation verwenden, nummerieren, Symbole definieren; aus dem Paper übernommene Gleichungen
mit Gleichungsnummer zitieren: \texttt{\textbackslash cite[Gl.~(3)]\{\PaperKey\}}.
\end{hinweise}
\begin{umfang}{Einlage 2 (Paper): Diskussion der Physik}{1}{3}
\platzhalter{Hauptaussage, Voraussetzungen, relevante Physik für \Schwerpunkt{} -- Stichpunkte aus dem
Abstract/der Einleitung des Papers:
(1) Drei Voraussetzungen für die Spin-Umkehr: unterschiedliche Hauptkrümmungen der Unterseite,
unterschiedliche horizontale Hauptträgheitsmomente, Fehlstellung (Winkel $\delta$) zwischen
Trägheits- und Krümmungshauptachsen;
(2) die Fehlstellung koppelt Drehen mit Nicken und Rollen -- asymmetrisch, mit entgegengesetztem
Vorzeichen des Drehmoments, daher Umkehr je nach Drehsinn unterschiedlich schnell;
(3) Annahmen: Rollen ohne Gleiten, keine Dissipation;
(4) Umkehrzeit $t_r$ (Garcia-Hubbard \cite{GarciaHubbard1988}) als Produkt aus Fehlstellungswinkel,
Differenz der inversen Trägheitsmomente, Differenz der Krümmungsradien und quadrierter
Schwingungsfrequenz -- Gleichung mit Nummer aus dem Paper zitieren;
(5) Simulationen: Formel gut bei kleiner Drehung und kleinen Schwingungen, bei schneller Drehung
keine Umkehr, Taumeln oder chaotisches Verhalten;
Bezug zum eigenen Schwerpunkt: Bewegungsgleichungen (Arbeitspaket 1.4) und eigene Simulation (1.6)
gegen die Ergebnisse des Papers prüfen}

% Beispiel für eine zitierte, nummerierte Gleichung (anpassen oder löschen):
% Nach \cite[Gl.~(1)]{\PaperKey} gilt
% \begin{equation}
%   t_r = \dots     % Garcia-Hubbard-Formel in der Form von Kondo & Nakanishi
%   \label{eq:paper-hauptgleichung}
% \end{equation}
% wobei ... ist. Aus Gl.~\eqref{eq:paper-hauptgleichung} folgt ...

\end{umfang}

\subsection{Reflexion}
\begin{hinweise}
In der Reflexion sollen Sie sich Gedanken machen über die Wirkung dieser Einlage und der Arbeit mit ihr auf
Ihren Lernfortschritt.
\begin{itemize}\itemsep0pt
\item In welchem Zusammenhang steht die Einlage mit Ihrem Schwerpunktthema?
\item Was haben Sie gelernt?
\item Wie hat Ihnen die Einlage beim Verständnis Ihres Projektes geholfen?
\item Wo gab es Schwierigkeiten bei der Beschäftigung mit der Einlage?
\item Hätte ein schon bestandenes Modul beim Verständnis helfen können? Welches? Warum oder warum nicht?
\item Was haben Sie durch diese Zusatzliteratur gelernt?
\item Was hätten Sie anders gemacht, wenn Sie die Einlage erstellt hätten?
\item Was sagen Ihre Kommilitoninnen oder Kommilitonen dazu, als Kommentar oder wenn Sie mit Ihnen über die Einlage diskutieren?
\item Haben Sie in dem Artikel der Einlage logische oder physikalische Fehler entdeckt?
\item Gab es bei der Beschäftigung mit der Einlage eine überraschende neue Erkenntnis?
\end{itemize}
\textbf{Die Reflexion für eine Einlage sollte mindestens eine Seite und maximal 2 Seiten lang sein. Die
Längenangaben sind ohne Titel, Bilder, Tabellen und Tabellen- oder Bildlegenden zu verstehen.}
\end{hinweise}
\begin{umfang}{Einlage 2 (Paper): Reflexion}{1}{2}
\platzhalter{Reflexion in Fließtext -- Leitfragen erkennbar beantworten}
\end{umfang}

%\subsection{Eigene Bewertungskriterien}
%...


\printbibliography[heading=bibintoc,title={Referenzen und Literaturverzeichnis}]

\appendix
\chapter*{Anh\"{a}nge}
\addcontentsline{toc}{chapter}{Anh\"{a}nge}
%%  ANHANG zu Einlage 2: das Paper selbst
%%  Kondo & Nakanishi (2017), Kopie aus Literatur/ -> anhaenge/paper_kondo_nakanishi_2017.pdf
%%  (Urheberrecht: Einbindung nur für die Abgabe/Bewertung.)
%%  Bei Wechsel auf Garcia & Hubbard (1988) Datei und Titel anpassen.

\anhangpdf{Wissenschaftliche Arbeit: Kondo und Nakanishi (2017) (Einlage 2)}{anh:paper}{anhaenge/paper_kondo_nakanishi_2017.pdf}


\end{document}
